\begin{figure}[tbp]
\centering
\begin{tikzpicture}[x=1cm,y=1cm,>=Latex,font=\small]
  \definecolor{fitblue}{RGB}{35,88,136}
  \definecolor{fitorange}{RGB}{168,83,31}
  \fill[fitblue!9] (-1.00,0)--(2.50,0)--(3.25,3.15)--(-0.25,3.15)--cycle;
  \draw[fitblue!50] (-1.00,0)--(-0.25,3.15) (2.50,0)--(3.25,3.15);
  \draw[fitblue!50] (-0.25,3.15)--(3.25,3.15);
  \fill[fitorange!12] (-1.225,0)--(0.525,0)--(-0.045,1.575)--(-1.795,1.575)--cycle;
  \draw[fitorange!65] (-1.225,0)--(-1.795,1.575) (0.525,0)--(-0.045,1.575);
  \draw[fitorange!65] (-1.795,1.575)--(-0.045,1.575);
  \draw[fitblue,densely dashed,line width=0.85pt] (0.75,0)--(1.50,3.15);
  \draw[fitorange,dash pattern=on 4pt off 2pt on 1pt off 2pt,line width=0.85pt]
      (-0.35,0)--(-0.92,1.575);
  \draw[->] (-2.65,0)--(3.65,0) node[right] {$x$};
  \draw[->] (-2.45,-0.02)--(-2.45,3.80) node[above] {$t$};
  \draw[black!35,dotted] (-2.45,3.15)--(1.50,3.15);
  \draw[black!35,dotted] (-2.45,1.575)--(-0.92,1.575);
  \node[left] at (-2.48,3.15) {$\delta_j$};
  \node[left] at (-2.48,1.575) {$\delta_k$};
  \node[left] at (-2.48,0) {$0$};
  \fill[fitblue] (1.50,3.15) circle (2.3pt);
  \node[above=5pt,fitblue] at (1.50,3.15) {$(x_j,\delta_j)$};
  \fill[fitorange] (-0.92,1.575) circle (2.3pt);
  \node[above left=3pt,fitorange] at (-0.92,1.575) {$(x_k,\delta_k)$};
  \fill (0,0) circle (2.8pt);
  \node[below=5pt] at (0,0) {$(x_{\rm end},0)$};
  \draw[fitblue] (0.75,-0.045)--(0.75,0.045);
  \draw[fitorange] (-0.35,-0.045)--(-0.35,0.045);
  \node[anchor=west,align=left,fitblue] (ja) at (4.00,2.80)
      {$x=x_j+p_j(t-\delta_j)$};
  \draw[fitblue!75] (ja.west)--(3.42,2.80)--(1.29,2.28);
  \node[anchor=west,align=left,fitorange] (ka) at (4.00,1.68)
      {$x=x_k+p_k(t-\delta_k)$};
  \draw[fitorange!75] (ka.west)--(3.42,1.68)--(-0.72,1.13);
  \node[anchor=west,align=left,text width=4.25cm] at (4.00,0.46)
      {$j,k\in\mathcal J$, $j<k$.\\[3pt]
       One endpoint lies in both strips.};
  \draw[fitblue,<->] (0.75,-0.72)--(2.50,-0.72)
      node[midway,below=2pt] {$K\delta_j$};
  \draw[fitblue!50] (0.75,-0.64)--(0.75,-0.82)
                        (2.50,-0.64)--(2.50,-0.82);
\end{tikzpicture}
\caption{Schematic localization at two selected scales, shown in one spatial
 direction with different horizontal and vertical units. The dashed axes
 belong to the two fitted frames; each shaded endpoint-localization strip has spatial half-width
 $K\delta_j$ or $K\delta_k$. Their time-zero sections both contain the
 same point $(x_{\rm end},0)$, expressing
 $|x_{\rm end}-x_j+p_j\delta_j|\le K\delta_j$.
 The conclusion holds for infinitely many selected indices
 $j\in\mathcal J$. The axes are separate affine fits, not a trajectory
 joining the marks.}
\label{fig:contact-localization}
\end{figure}
